\chapter{Creio na vida eterna}

O décimo segundo artigo do Credo confessa a fé na vida eterna, apontando para o destino último e definitivo da humanidade. Para o cristão, a morte não é um fim absoluto, mas uma passagem e um encontro.

\section{Introdução (1020)}

O cristão que une a sua própria morte à morte de Jesus vê nela um passo em direção a Ele e a entrada na vida eterna. Quando a Igreja, pela última vez, pronuncia as palavras de perdão e absolvição sobre o cristão moribundo, assinalando-o com uma unção fortificante e dando-lhe Cristo em viático como alimento para a viagem, ela fala com uma doce segurança.

\section{O Juízo Particular (1021--1022)}

A morte põe termo à vida humana como o tempo aberto para aceitar ou rejeitar a graça divina manifestada em Cristo. O Novo Testamento fala do juízo primariamente no seu aspecto de encontro final com Cristo na Sua segunda vinda, mas afirma repetidamente que cada um será recompensado imediatamente após a morte, de acordo com as suas obras e a sua fé. A parábola do pobre Lázaro e as palavras de Cristo na cruz ao bom ladrão, assim como outros textos, falam de um destino final da alma, que pode ser diferente para uns e para outros.

Cada homem recebe na sua alma imortal a sua retribuição eterna no momento exato da sua morte, num juízo particular que refere a sua vida a Cristo: quer a entrada na bem-aventurança do Céu --- através de uma purificação ou imediatamente --- quer a condenação imediata e eterna.

\section{O Céu (1023--1029)}

Aqueles que morrem na graça e na amizade de Deus, e que estão perfeitamente purificados, vivem para sempre com Cristo. Eles são para sempre semelhantes a Deus, pois ``vêem-no tal como Ele é'' (1 Jo 3,2), face a face.

Esta vida perfeita com a Santíssima Trindade --- esta comunhão de vida e de amor com a Trindade, com a Virgem Maria, os anjos e todos os bem-aventurados --- é chamada de ``Céu''. O Céu é o fim último e a realização dos anseios humanos mais profundos, o estado de suprema e definitiva felicidade.

Viver no Céu é ``estar com Cristo''. Os eleitos vivem ``em Cristo'', mas retêm, ou melhor, encontram a sua verdadeira identidade, o seu próprio nome. Pela sua Morte e Ressurreição, Jesus Cristo abriu-nos o céu. A vida dos bem-aventurados consiste na posse em plenitude dos frutos da redenção realizada por Cristo. Ele associa à sua glorificação celeste aqueles que acreditaram n'Ele e permaneceram fiéis à sua vontade. O céu é a comunidade bem-aventurada de todos os que estão perfeitamente incorporados em Cristo.

Este mistério de comunhão bem-aventurada com Deus e com todos os que estão em Cristo ultrapassa toda a compreensão e toda a representação. A Escritura fala-nos dele através de imagens, como vida, luz, paz, banquete de núpcias, vinho do Reino, casa do Pai, a Jerusalém celeste, o paraíso: ``O que os olhos não viram, os ouvidos não ouviram, nem o coração do homem imaginou, o que Deus preparou para aqueles que O amam'' (1 Cor 2, 9).

Em virtude da sua transcendência, Deus não pode ser visto tal como é a não ser quando Ele próprio abre o seu mistério à contemplação imediata do homem e lhe dá a capacidade para o fazer. A contemplação de Deus na sua glória celeste é chamada pela Igreja de ``visão beatífica''.

Na glória do céu, os bem-aventurados continuam a cumprir com alegria a vontade de Deus em relação aos outros homens e a toda a criação. Já reinam com Cristo; com Ele `` reinarão pelos séculos dos séculos'' (Ap 22, 5).

\section{A purificação final ou Purgatório (1030--1032)}

Os que morrem na graça e na amizade de Deus, mas não de todo purificados, embora seguros da sua salvação eterna, sofrem depois da morte uma purificação, a fim de obterem a santidade necessária para entrar na alegria do céu.

A Igreja chama Purgatório a esta purificação final dos eleitos, que é absolutamente distinta do castigo dos condenados. A doutrina da fé relativamente ao Purgatório foi formulada pela Igreja sobretudo nos concílios de Florença e de Trento. A Tradição da Igreja, referindo-se a certos textos da Escritura, fala de um fogo purificador.

Esta doutrina apoia-se também na prática da oração pelos defuntos, de que já fala a Sagrada Escritura: ``Por isso, Judas Macabeu pediu um sacrifício expiatório para que os mortos fossem livres das suas faltas'' (2Mac 12, 46). Desde os primeiros tempos, a Igreja honrou a memória dos defuntos, oferecendo sufrágios em seu favor, particularmente o Sacrifício eucarístico para que, purificados, possam chegar à visão beatífica de Deus. A Igreja recomenda também a esmola, as indulgências e as obras de penitência a favor dos defuntos.

\section{O Inferno (1033--1037)}

Não podemos estar em união com Deus se não escolhermos livremente amá-Lo. Mas não podemos amar a Deus se pecarmos gravemente contra Ele, contra o nosso próximo ou contra nós mesmos. Morrer em pecado mortal sem arrependimento e sem dar acolhimento ao amor misericordioso de Deus, significa permanecer separado d'Ele para sempre, por nossa própria livre escolha. E é este estado de auto-exclusão definitiva da comunhão com Deus e com os bem-aventurados que se designa pela palavra ``Inferno''.

Jesus fala muitas vezes da ``geena'' do ``fogo que não se apaga'', reservada aos que recusam, até ao fim da vida, acreditar e converter-se. A doutrina da Igreja afirma a existência do Inferno e a sua eternidade. As almas dos que morrem em estado de pecado mortal descem imediatamente, após a morte, aos infernos, onde sofrem as penas do Inferno, ``o fogo eterno''. A principal pena do inferno consiste na separação eterna de Deus, o único em Quem o homem pode ter a vida e a felicidade para que foi criado e a que aspira.

As afirmações da Sagrada Escritura e os ensinamentos da Igreja a respeito do Inferno são um apelo ao sentido de responsabilidade com que o homem deve usar da sua liberdade, tendo em vista o destino eterno. Constituem, ao mesmo tempo, um apelo urgente à conversão.

Deus não predestina ninguém para o Inferno. Para ter semelhante destino, é preciso haver uma aversão voluntária a Deus (pecado mortal) e persistir nela até ao fim. Na liturgia eucarística e nas orações quotidianas dos seus fiéis, a Igreja implora a misericórdia de Deus, ``que não quer que ninguém pereça, mas que todos se convertam'' (2Pe 3,9).

\section{O Juízo final (1038--1041)}

A ressurreição de todos os mortos, ``justos e pecadores'' (At 24,15), há-de preceder o Juízo final. Será ``a hora em que todos os que estão nos túmulos hão-de ouvir a sua voz e sairão: os que tiverem praticado o bem, para uma ressurreição de vida, e os que tiverem praticado o mal, para uma ressurreição de condenação'' (Jo 5,28-29). Então Cristo virá ``na sua glória, com todos os seus anjos''.

É perante Cristo, que é a Verdade, que será definitivamente posta a descoberto a verdade da relação de cada homem com Deus. O Juízo final revelará, até às suas últimas consequências, o que cada um tiver feito ou deixado de fazer de bem durante a sua vida terrena.

O Juízo final terá lugar quando acontecer a vinda gloriosa de Cristo. Só o Pai sabe o dia e a hora, só Ele decide sobre a sua vinda. O Juízo final revelará como a justiça de Deus triunfa de todas as injustiças cometidas pelas suas criaturas e como o seu amor é mais forte do que a morte.

A mensagem do Juízo final é um apelo à conversão, enquanto Deus dá ainda aos homens ``o tempo favorável, o tempo da salvação'' (2Cor 6,2).

\section{A esperança dos novos céus e da nova terra (1042--1050)}

No fim dos tempos, o Reino de Deus chegará à sua plenitude. Depois do Juízo final, os justos reinarão para sempre com Cristo, glorificados em corpo e alma, e o próprio universo será renovado.

A esta misteriosa renovação, que há-de transformar a humanidade e o mundo, a Sagrada Escritura chama ``os novos céus e a nova terra'' (2Pe 3,13). Será a realização definitiva do desígnio divino de ``reunir sob a chefia de Cristo todas as coisas que há nos céus e na terra'' (Ef 1,10).

Neste ``mundo novo'', a Jerusalém celeste, Deus terá a sua morada entre os homens. Para o homem, esta consumação será a realização final da unidade do género humano, querida por Deus desde a criação e da qual a Igreja peregrina era ``como que o sacramento''. Os que estiverem unidos a Cristo formarão a comunidade dos resgatados, a ``Cidade santa de Deus'' (Ap 21,2), a ``Esposa do Cordeiro'' (Ap 21,9).

Quanto ao cosmos, a Revelação afirma a profunda comunidade de destino entre o mundo material e o homem. Assim, pois, também o universo visível está destinado a ser transformado.

Ignoramos o tempo em que a terra e a humanidade atingirão a sua plenitude, e também não sabemos como é que o universo será transformado. A expectativa da nova terra não deve, porém, enfraquecer, mas antes activar a solicitude em ordem a desenvolver esta terra onde cresce o corpo da nova família humana.

Pois todos os bens da dignidade humana, da comunhão fraterna e da liberdade, ou seja, todos os frutos da natureza e da empresa humana são destinados à ``nova terra'' e ``novo céu''.

\clearpage

\section{O Símbolo de Fé}

\begin{tabular}{|c|p{4.5cm}|p{5.5cm}|}
    \hline
        Artigo & Símbolo dos Apóstolos & Símbolo Niceno-Constantinopolitano \\ \hline
        1º & Creio em Deus Pai todo-póderoso,
        criador do céu e da terra.
 & Creio em um só Deus, Pai todo-poderoso,
 criador do céu e da terra,
 de todas as coisas visíveis e invisíveis.
 \\ \hline
        2º & E em Jesus Cristo, seu único Filho, nosso Senhor,        
 & Creio em um só Senhor, Jesus Cristo, Filho Unigênito de Deus,
 Nascido do Pai antes de todos os séculos:
 Deus de Deus,
 luz da luz,
 Deus verdadeiro de Deus verdadeiro,
 gerado, não criado,
 consubstancial ao Pai.
 Por ele todas as coisas foram feitas.
 \\ \hline
        3º & Que foi concebido pelo Espírito Santo, nasceu da Virgem Maria; & E por nós, homens, e para nossa salvação,
        Desceu dos céus: e se encarnou pelo Espírito Santo, no seio da virgem Maria, e se fez homem; \\ \hline
        4º & padeceu sob Pôncio Pilatos,
        foi crucificado, morto e sepultado.
 & Também por nós foi crucificado sob Pôncio Pilatos;
 Padeceu e foi sepultado.
 \\ \hline
        5º & Desceu à mansão dos mortos; ressuscitou ao terceiro dia, & e Ressuscitou ao terceiro dia,
        conforme as Escrituras,
 \\ \hline
        6º & subiu aos céus;
        está sentado à direita de Deus Pai todo-poderoso,
 & e subiu aos céus,
 onde está sentado à direita do Pai.
 \\ \hline
        7º & donde há de vir a julgar os vivos e os mortos. & E de novo há de vir, em sua glória,
        Para julgar os vivos e os mortos;
        e o seu reino não terá fim.
 \\ \hline
        8º & Creio no Espírito Santo; & Creio no Espírito Santo,
        Senhor que dá a vida,
        e procede do Pai e do Filho;
        e com o Pai e o Filho é adorado e glorificado:
        ele que falou pelos profetas.
 \\ \hline
        9º & na Santa Igreja Católica; na comunhão dos santos; & Creio na Igreja,
        una, santa, católica e apostólica.
 \\ \hline
        10º & na remissão dos pecados; & Professo um só batismo 
        para remissão dos pecados.
 \\ \hline
        11º & na ressurreição da carne & E espero a ressurreição dos mortos \\ \hline
        12º & e na vida eterna. Amém. & e a vida do mundo que há de vir. Amém. \\ \hline
    \end{tabular}
